% ──────────────────────────────────────────────────────────────────────────────
%  Informe Lab 4 — Análisis de Datos Geoespaciales
%  CC3084 Data Science · Universidad del Valle de Guatemala · Semestre II 2026
%  Autores: Roberto Barreda (23354), Javier España (23361), Ángel Esquit (23221)
%
%  INSTRUCCIONES PARA COMPILAR:
%    pdflatex informe_lab04.tex
%    pdflatex informe_lab04.tex   (segunda pasada para referencias)
%
%  Figuras esperadas en docs/figs/ (generadas al exportar desde el notebook):
%    mapa_lagos.png            - mapa folium de ubicación (captura)
%    mapa_indices_*.png        - paneles color real / mascara / NDCI / Chl-a
%    serie_temporal.png        - gráfico de línea Chl-a por lago
%    comparacion_chla_ndwi.png - gráfico log + boxplot + % área alta
%    correlacion.png           - scatter Chl-a vs NDVI / NDWI
%    comparacion_lagos.png     - comparación intensidad y frecuencia
%    extension_espacial.png    - % lago sobre 30 mg/m³ por fecha
%    zonas_persistentes.png    - mapa de Chl-a promedio por píxel
%    boxplots_fechas.png       - distribución estadística entre fechas
%    estacionalidad.png        - boxplot seca vs lluviosa
% ──────────────────────────────────────────────────────────────────────────────

\documentclass[12pt, a4paper]{article}

% ── Paquetes ──────────────────────────────────────────────────────────────────
\usepackage[utf8]{inputenc}
\usepackage[T1]{fontenc}
\usepackage[spanish]{babel}
\usepackage{lmodern}
\usepackage{geometry}
\usepackage{graphicx}
\usepackage{caption}
\usepackage{subcaption}
\usepackage{float}
\usepackage{booktabs}
\usepackage{array}
\usepackage{xcolor}
\usepackage{hyperref}
\usepackage{amsmath}
\usepackage{enumitem}
\usepackage{parskip}
\usepackage{titlesec}
\usepackage{fancyhdr}
\usepackage{microtype}
\usepackage{setspace}

\geometry{
  top=2.5cm, bottom=2.5cm,
  left=3cm,  right=2.5cm,
}

% ── Colores institucionales ───────────────────────────────────────────────────
\definecolor{uvgblue}{RGB}{0, 56, 101}
\definecolor{atitlanblue}{RGB}{31, 119, 180}
\definecolor{amatitlanred}{RGB}{214, 39, 40}
\definecolor{alertred}{RGB}{192, 0, 0}
\definecolor{lightgray}{gray}{0.93}

% ── Estilo de secciones ───────────────────────────────────────────────────────
\titleformat{\section}{\large\bfseries\color{uvgblue}}{}{0em}{}[\titlerule]
\titleformat{\subsection}{\normalsize\bfseries\color{uvgblue}}{}{0em}{}
\titleformat{\subsubsection}{\normalsize\itshape}{}{0em}{}

% ── Encabezado / pie de página ────────────────────────────────────────────────
\pagestyle{fancy}
\fancyhf{}
\fancyhead[L]{\small\color{uvgblue}CC3084 Data Science · Lab 4 · UVG 2026}
\fancyhead[R]{\small\color{uvgblue}\thepage}
\renewcommand{\headrulewidth}{0.4pt}

% ── Hipervínculos ─────────────────────────────────────────────────────────────
\hypersetup{
  colorlinks=true,
  linkcolor=uvgblue,
  urlcolor=uvgblue,
  citecolor=uvgblue,
  pdftitle={Lab 4 — Análisis de Datos Geoespaciales},
  pdfauthor={Barreda · España · Esquit},
}

% ── Utilidades ────────────────────────────────────────────────────────────────
\newcommand{\lago}[1]{\textbf{#1}}
\newcommand{\indice}[1]{\texttt{#1}}
\newcommand{\mgm}{\,mg/m$^3$}
\newcommand{\kmsq}{\,km$^2$}
\newcommand{\nota}[1]{\vspace{2pt}\noindent\colorbox{lightgray}{\parbox{\dimexpr\linewidth-2\fboxsep\relax}{\small\textit{#1}}}\vspace{4pt}}

% ─────────────────────────────────────────────────────────────────────────────
\begin{document}

% ── Portada ───────────────────────────────────────────────────────────────────
\begin{titlepage}
  \centering
  \vspace*{1cm}
  {\color{uvgblue}\rule{\linewidth}{2pt}}\vspace{8pt}

  {\LARGE\bfseries\color{uvgblue}
    Monitoreo de Floraciones de Cianobacterias\\[6pt]
    en los Lagos Atitlán y Amatitlán\\[4pt]
    mediante Imágenes Sentinel-2
  }\vspace{10pt}

  {\color{uvgblue}\rule{\linewidth}{1pt}}\vspace{14pt}

  {\large Laboratorio 4 — Análisis de Datos Geoespaciales}\vspace{4pt}

  {\normalsize CC3084 Data Science}\vspace{2pt}

  {\normalsize Universidad del Valle de Guatemala}\vspace{2pt}

  {\normalsize Facultad de Ingeniería · Semestre II 2026}\vspace{20pt}

  \begin{tabular}{ll}
    \textbf{Autores:} & Roberto Barreda \quad (23354) \\
                      & Javier España    \quad (23361) \\
                      & Ángel Esquit   \quad (23221)
  \end{tabular}\vspace{20pt}

  {\normalsize Fecha de entrega: 16 de agosto de 2026}\vspace{10pt}

  {\color{uvgblue}\rule{\linewidth}{2pt}}

  \vfill
  \nota{Este informe está dirigido a profesionales ambientalistas. Los conceptos técnicos de teledetección y programación se explican en lenguaje accesible.}
\end{titlepage}

% ── Resumen ejecutivo ─────────────────────────────────────────────────────────
\section*{Resumen}
\addcontentsline{toc}{section}{Resumen}

Este informe presenta los resultados del monitoreo satelital de floraciones de cianobacterias
en los lagos \lago{Atitlán} y \lago{Amatitlán} de Guatemala, usando 11 imágenes
Sentinel-2 por lago tomadas entre enero de 2025 y julio de 2026.

Los resultados muestran una diferencia de régimen, no de grado, entre ambos cuerpos de agua:

\begin{itemize}[leftmargin=1.5em]
  \item \lago{Amatitlán} presenta eutrofia crónica (18.4\mgm{} en promedio; máximo 60.2\mgm{}
        el 2026-06-19, cuando el 76.6\,\% del lago superó el umbral de floración densa).
  \item \lago{Atitlán} permanece en buen estado a escala de lago (3.5\mgm{} en promedio),
        pero muestra focos puntuales intensos en bahías cercanas a centros poblados.
\end{itemize}

\tableofcontents
\newpage

% ═══════════════════════════════════════════════════════════════════════════════
\section{Introducción y contexto ambiental}
% ═══════════════════════════════════════════════════════════════════════════════

Las \textbf{cianobacterias} son microorganismos fotosintéticos que proliferan cuando el agua
reúne tres condiciones: temperatura elevada, abundancia de nutrientes (nitrógeno y fósforo) y
estancamiento. Cuando forman ``floraciones'' pueden producir toxinas que afectan la salud humana
y animal, matan peces, generan olores desagradables y devalúan zonas turísticas y de pesca.

Los lagos \lago{Atitlán} y \lago{Amatitlán} son casos emblemáticos en Guatemala:

\begin{description}[leftmargin=1.5em, font=\bfseries]
  \item[Lago Atitlán] Se ubica en el altiplano occidental, dentro de una caldera volcánica de
        más de 300\,m de profundidad. Tiene una superficie de $\sim$130\kmsq{} y es uno de los
        destinos turísticos más visitados del país. Su primera floración masiva documentada ocurrió
        en 2009 y generó alarma internacional.
  \item[Lago Amatitlán] Es un lago mucho más pequeño ($\sim$15.2\kmsq) y somero ($\sim$33\,m),
        ubicado a menos de 25\,km al sur de la Ciudad de Guatemala. Recibe el río Villalobos,
        que drena buena parte del área metropolitana (unos 3 millones de habitantes). Su estado
        de deterioro es crónico y bien documentado desde hace décadas.
\end{description}

El monitoreo físico de estos lagos es costoso y limitado en cobertura espacial. Los satélites
\textbf{Sentinel-2} del programa europeo Copernicus ofrecen imágenes multiespectrales de 20\,m
de resolución que permiten calcular índices de calidad del agua de forma sistemática y gratuita.

% ═══════════════════════════════════════════════════════════════════════════════
\section{Metodología}
% ═══════════════════════════════════════════════════════════════════════════════

\subsection{Datos utilizados}

Se usaron las 11 fechas oficiales por lago proporcionadas por el enunciado, todas con
nubosidad baja ($<$15\,\%). Para cada fecha y lago se descargaron únicamente las bandas
espectrales necesarias para los índices solicitados (9 bandas por imagen, resolución 20\,m),
reduciendo el volumen de descarga respecto a una escena completa.

\begin{table}[H]
\centering
\caption{Fechas de adquisición por lago}
\begin{tabular}{ll}
\toprule
\textbf{Amatitlán} & \textbf{Atitlán} \\
\midrule
2025-01-28, 2025-04-15, 2025-04-28 & 2025-01-18, 2025-04-13, 2025-05-13 \\
2025-11-24, 2026-01-08, 2026-02-02 & 2025-07-17, 2025-11-21, 2025-12-29 \\
2026-02-07$^*$, 2026-03-29, 2026-04-13 & 2026-02-12, 2026-03-24, 2026-04-13 \\
2026-04-28, 2026-06-19 & 2026-04-28, 2026-07-22 \\
\bottomrule
\multicolumn{2}{l}{\small $^*$ Cobertura válida parcial ($\sim$57\,\%); incluida porque el lago está completamente cubierto.}
\end{tabular}
\end{table}

\subsection{Índices calculados}

\subsubsection{Índice de cianobacteria (Chl-a)}

Se reprodujo en Python el \emph{Cyano Detection Script} (\texttt{cyanobacteria\_chla\_ndci\_l1c})
publicado en \href{https://custom-scripts.sentinel-hub.com}{custom-scripts.sentinel-hub.com}
(CyanoLakes, Kravitz \& Matthews 2020). El índice NDCI (Normalized Difference Chlorophyll Index)
relaciona la banda del ``borde rojo'' con la del rojo:

\[
  \text{NDCI} = \frac{B05 - B04}{B05 + B04}
\]

La conversión a clorofila-a en mg/m$^3$ usa el polinomio cúbico calibrado por los autores:
\[
  \text{Chl-a} = 826.57\,\text{NDCI}^3 - 176.43\,\text{NDCI}^2 + 19\,\text{NDCI} + 4.071
\]

Los valores se recortan al rango de calibración declarado (0--500\mgm).

\subsubsection{NDVI y NDWI}

\[
  \text{NDVI} = \frac{B08 - B04}{B08 + B04} \qquad
  \text{NDWI} = \frac{B03 - B08}{B03 + B08}
\]

El NDVI mide biomasa vegetal; el NDWI, la presencia de agua superficial (valores positivos
indican agua). Ambos se calcularon sobre los píxeles clasificados como agua.

\subsection{Máscara de agua y huella del lago}

Para evitar confundir suelo húmedo o sombras urbanas con agua se aplicó la función
\texttt{wbi()} del script, que combina seis criterios espectrales (MNDWI, NDWI, AWEI\textsubscript{sh},
AWEI\textsubscript{nsh}, NDVI, NDWI\textsubscript{hojas}).

Dado que la máscara varía ligeramente entre fechas (la floración densa cambia el espectro del
agua y puede excluir píxeles con cianobacterias), se construyó una \textbf{huella estable}:
los píxeles clasificados como agua en al menos el 30\,\% de las fechas, cerrados morfológicamente
y reducidos al cuerpo principal. Todas las mediciones temporales usan esta huella fija como
denominador.

\subsection{Control de calidad}

Las fechas en que menos del 70\,\% de la huella tiene datos válidos se marcan como
\emph{no confiables} y se excluyen del análisis temporal. Solo Atitlán 2025-01-18 fue
descartada: la corrección atmosférica sobre el agua ultraoscura de la caldera generó
reflectancias negativas en el rojo y el borde rojo, haciendo que el NDCI fuese
indeterminado para el 76\,\% del lago. Sin ese filtro la imagen habría reportado una
media de $\sim$30\,000\mgm, arruinando la serie temporal completa.

% ═══════════════════════════════════════════════════════════════════════════════
\section{Resultados}
% ═══════════════════════════════════════════════════════════════════════════════

% ── 3.1 Ubicación y área de estudio ──────────────────────────────────────────
\subsection{Ubicación de los lagos}

El mapa interactivo generado con \texttt{folium} (Figura~\ref{fig:mapa_lagos}) muestra la
posición geográfica de ambos lagos y los recuadros de descarga de cada imagen.

\begin{figure}[H]
  \centering
  \fbox{\begin{minipage}{0.85\linewidth}\centering\vspace{1.2cm}
    \textit{[Captura del mapa folium interactivo.\\
    Ejecutar la celda de ubicación del notebook y guardar como PNG en docs/figs/mapa\_lagos.png]}
    \vspace{1.2cm}
  \end{minipage}}
  \caption{Posición de los lagos Atitlán y Amatitlán en Guatemala, con sus \emph{bounding boxes} de descarga de imágenes Sentinel-2.}
  \label{fig:mapa_lagos}
\end{figure}

% ── 3.2 Mapas por fecha ───────────────────────────────────────────────────────
\subsection{Índice de cianobacteria por fecha (Ejercicio 3)}

Para cada lago se generaron mapas de cuatro paneles (color real, máscara de agua, NDCI y
Chl-a con la rampa de color oficial). La Figura~\ref{fig:mapas_indices} muestra la fecha
más limpia y la del pico de floración de cada lago.

\begin{figure}[H]
  \centering
  \begin{subfigure}{\linewidth}
    \fbox{\begin{minipage}{\linewidth}\centering\vspace{1cm}
      \textit{[mapa\_indices\_amatitlan\_limpia.png — fecha más limpia de Amatitlán]}
      \vspace{1cm}
    \end{minipage}}
    \caption{Amatitlán — fecha con menor Chl-a (agua relativamente limpia).}
  \end{subfigure}

  \vspace{6pt}

  \begin{subfigure}{\linewidth}
    \fbox{\begin{minipage}{\linewidth}\centering\vspace{1cm}
      \textit{[mapa\_indices\_amatitlan\_pico.png — 2026-06-19, pico de floración]}
      \vspace{1cm}
    \end{minipage}}
    \caption{Amatitlán — 2026-06-19: floración generalizada en el 76.6\,\% del lago.}
  \end{subfigure}

  \caption{Paneles de color real / máscara de agua / NDCI / Chl-a para Lago Amatitlán.
           Azul = agua limpia; verde-amarillo = biomasa moderada; naranja-rojo = floración densa.}
  \label{fig:mapas_indices}
\end{figure}

% ── 3.3 Análisis temporal ─────────────────────────────────────────────────────
\subsection{Análisis temporal — evolución de la clorofila-a (Ejercicio 4)}

\subsubsection{Lago Amatitlán}

Las 11 fechas dan una media de \textbf{18.4\mgm} (desviación estándar: 15.2\mgm{}),
entre 5.9\mgm{} (2026-02-02) y \textbf{60.2\mgm{} el 2026-06-19} — una variación
de 10\,× dentro del mismo período. El nivel de fondo \emph{nunca} baja de 5.9\mgm:
no hay ninguna fecha en que Amatitlán se pueda considerar limpio. En promedio, el
\textbf{15.6\,\%} de su superficie supera los 30\mgm{}.

El criterio de pico (media + 1$\sigma$, umbral = 33.6\mgm) señala una única fecha
crítica: \textbf{2026-06-19}, con 60.2\mgm{} de promedio y el 76.6\,\% del lago en
estado de floración densa. Le siguen, sin alcanzar el umbral, 2026-04-28 (25.5\mgm) y
2026-01-08 (24.2\mgm).

\subsubsection{Lago Atitlán}

Las 10 fechas confiables muestran una Chl-a de \textbf{3.48 ± 0.46\mgm{}}, variando
entre 2.49 y 3.96\mgm{}: una variación total de 1.6\,×. En promedio, solo el
\textbf{0.22\,\%} de la superficie supera los 30\mgm. A escala de lago,
\emph{Atitlán no presenta floraciones} en el período estudiado.

Los dos ``picos'' estadísticos detectados (2025-07-17 y 2026-07-22) superan el umbral
por centésimas y están dentro del ruido de la serie. Lo que sí aparecen son focos
\textbf{pequeños e intensos} en las bahías: en 2025-11-21 el máximo por píxel llegó
al tope de la escala (500\mgm) con el 1.0\,\% del lago ($\sim$1.2\kmsq) afectado.

\begin{figure}[H]
  \centering
  \fbox{\begin{minipage}{0.9\linewidth}\centering\vspace{2cm}
    \textit{[serie\_temporal.png — gráfico de línea Chl-a por lago con umbrales de pico]}
    \vspace{2cm}
  \end{minipage}}
  \caption{Evolución temporal de la clorofila-a media por lago y fecha. Los círculos vacíos marcan los picos de floración (media + 1$\sigma$). Las cruces grises indican fechas descartadas por cobertura insuficiente. La línea discontinua es el umbral de pico de cada lago.}
  \label{fig:serie_temporal}
\end{figure}

% ── 3.4 Análisis espacial ────────────────────────────────────────────────────
\subsection{Análisis espacial — distribución dentro del lago (Ejercicio 5)}

\subsubsection{Amatitlán}

El foco de mayor Chl-a se concentra sistemáticamente en el sector \textbf{norte y noreste}
del lago, cerca de la desembocadura del río Villalobos. En el pico del 2026-06-19 la
floración es prácticamente uniforme en todo el espejo de agua (razón mediana/media = 0.98).
En fechas de floraciones en parches (p.ej., 2026-04-28, razón 0.41) el mapa muestra manchas
localizadas junto a las entradas de afluentes.

\subsubsection{Atitlán}

El mapa de color es predominantemente azul ($<$5\mgm), pero aparecen focos puntuales e
intensos en las bahías de Santiago Atitlán, San Pedro La Laguna y Panajachel — municipios
con mayor actividad humana y descarga de aguas residuales. El centro del lago permanece limpio.

\begin{figure}[H]
  \centering
  \fbox{\begin{minipage}{0.9\linewidth}\centering\vspace{2cm}
    \textit{[comparacion\_cuadricula\_amatitlan.png — grid comparativo de todas las fechas]}
    \vspace{2cm}
  \end{minipage}}
  \caption{Cuadrícula comparativa de clorofila-a en todas las fechas de Amatitlán (bordes rojos = pico de floración). La misma huella fija se usa en todas las fechas para que la comparación sea justa.}
  \label{fig:grid_amatitlan}
\end{figure}

% ── 3.5 Correlación de índices ───────────────────────────────────────────────
\subsection{Correlación de índices NDVI y NDWI con la cianobacteria (Ejercicio 6)}

\begin{table}[H]
\centering
\caption{Coeficientes de correlación de Pearson (fechas confiables)}
\begin{tabular}{lcccc}
\toprule
\textbf{Lago} & \textbf{r(NDVI vs Chl-a)} & \textbf{p-valor} & \textbf{r(NDWI vs Chl-a)} & \textbf{p-valor} \\
\midrule
Amatitlán & positiva moderada–alta & $< 0.05$ & negativa moderada–alta & $< 0.05$ \\
Atitlán   & muy débil              & $> 0.05$ & muy débil              & $> 0.05$ \\
\bottomrule
\end{tabular}
\end{table}

\textbf{Interpretación.} En \lago{Amatitlán} la correlación positiva NDVI--Chl-a confirma que
las algas elevan la reflectancia en el infrarrojo respecto al rojo. La correlación negativa
NDWI--Chl-a es esperada: el agua cargada de pigmentos absorbe diferente al agua limpia. Los tres
índices apuntan en la misma dirección y cualquiera sirve como indicador de deterioro.

En \lago{Atitlán} las correlaciones son débiles porque la Chl-a varía muy poco a escala de lago
($\sigma = 0.46\mgm$): el ruido de medición enmascara la señal. A escala de píxel en las bahías
afectadas, la relación probablemente sí se mantiene, pero el promedio del lago la diluye.

% ═══════════════════════════════════════════════════════════════════════════════
\section{Comparación entre los lagos (Ejercicio 7)}
% ═══════════════════════════════════════════════════════════════════════════════

\begin{table}[H]
\centering
\caption{Comparación de indicadores entre ambos lagos (período enero 2025–julio 2026)}
\begin{tabular}{lccc}
\toprule
\textbf{Indicador} & \textbf{Amatitlán} & \textbf{Atitlán} & \textbf{Razón} \\
\midrule
Chl-a media período      & 18.4\mgm    & 3.5\mgm    & \textbf{5.3×} \\
Chl-a máxima por fecha   & 60.2\mgm    & 4.0\mgm    & \textbf{15.1×} \\
Superficie sobre 30\mgm  & 15.6\,\%    & 0.22\,\%   & \textbf{70×} \\
Variación entre fechas   & 10.2×       & 1.6×       & — \\
N.º de picos detectados  & 1           & 2*         & — \\
Superficie huella        & 13.8\kmsq   & 120.8\kmsq & 0.11× \\
\bottomrule
\multicolumn{4}{l}{\small * Picos estadísticos dentro del ruido, no eventos ecológicos de escala de lago.}
\end{tabular}
\end{table}

\subsection{Causas de las diferencias}

\begin{table}[H]
\centering
\caption{Factores que explican la diferencia entre lagos}
\begin{tabular}{lll}
\toprule
\textbf{Factor} & \textbf{Amatitlán} & \textbf{Atitlán} \\
\midrule
Profundidad media   & $\sim$33 m           & $>$188 m \\
Principal afluente  & Río Villalobos       & Ríos menores y lluvia \\
Presión urbana      & Muy alta ($\sim$3 M hab.) & Moderada \\
Régimen trófico     & Hipereutrófico       & Oligo-mesotrófico \\
\bottomrule
\end{tabular}
\end{table}

\lago{Amatitlán} recibe el drenaje de gran parte del área urbana de la Ciudad de Guatemala
a través del río Villalobos: una carga de nitrógeno y fósforo que su pequeño volumen no puede
diluir. El resultado es eutrofia crónica con un nivel de fondo que nunca baja de 5.9\mgm{} y
un pico de 60.2\mgm{} al inicio de la época lluviosa.

\lago{Atitlán} cuenta con la ventaja de su enorme profundidad (más de 300\,m en el centro de
la caldera): el volumen es tan grande que los nutrientes se diluyen antes de acumularse. Sin
embargo, los focos puntuales en bahías con descarga de aguas residuales muestran que la presión
antrópica ya es visible localmente. Si esos aportes se intensifican con el crecimiento turístico,
el deterioro podría acelerarse.

% ═══════════════════════════════════════════════════════════════════════════════
\section{Análisis exploratorio adicional (Ejercicio 8)}
% ═══════════════════════════════════════════════════════════════════════════════

\subsection{Extensión espacial de la floración}

\begin{figure}[H]
  \centering
  \fbox{\begin{minipage}{0.9\linewidth}\centering\vspace{1.5cm}
    \textit{[extension\_espacial.png — \% del lago sobre 30 mg/m³ por fecha]}
    \vspace{1.5cm}
  \end{minipage}}
  \caption{Porcentaje del espejo de agua de cada lago con valores de Chl-a superiores a 30\mgm{} por fecha. La línea discontinua marca el promedio del período.}
  \label{fig:extension}
\end{figure}

El 2026-06-19 Amatitlán registra un 76.6\,\% de su superficie sobre el umbral de floración densa:
no se trata de un foco localizado sino de una floración de lago completo. En Atitlán ninguna fecha
supera el 1.5\,\% sobre ese umbral.

\subsection{Zonas persistentes de acumulación}

El mapa de clorofila-a promediada por píxel a lo largo de todas las fechas disponibles permite
identificar las zonas donde la biomasa algal se acumula de forma recurrente (Figura~\ref{fig:zonas}).
En Amatitlán el sector norte concentra los valores más altos. En Atitlán los focos persistentes
corresponden a las bahías con mayor actividad humana.

\begin{figure}[H]
  \centering
  \fbox{\begin{minipage}{0.9\linewidth}\centering\vspace{2cm}
    \textit{[zonas\_persistentes.png — mapa de Chl-a promedio por píxel]}
    \vspace{2cm}
  \end{minipage}}
  \caption{Mapa de clorofila-a promedio por píxel (media de todas las fechas con datos disponibles). Los tonos más oscuros/rojos indican zonas con acumulación recurrente.}
  \label{fig:zonas}
\end{figure}

\subsection{Distribución estadística de la clorofila-a}

La distribución de Chl-a en Amatitlán es asimétrica a la derecha en casi todas las fechas
(muchos píxeles con valores moderados y una cola hacia valores altos). En el pico del 2026-06-19
la distribución se uniformiza (razón mediana/media = 0.98), señal de floración pareja en todo
el lago. En Atitlán la distribución es siempre estrecha y centrada cerca de 3.5\mgm{}, con colas
extremas en algunas fechas asociadas a los focos de bahía.

\subsection{Patrón estacional}

Separando las fechas en época seca (noviembre–abril) y lluviosa (mayo–octubre), Amatitlán
promedia 14.3\mgm{} en seca frente a 60.2\mgm{} en lluviosa, y Atitlán 3.34 frente a 3.81\mgm{}.
La dirección es la esperada: las lluvias lavan suelos agrícolas y calles, inyectando nutrientes
a los lagos.

\nota{Limitación estadística: Amatitlán solo tiene 1 fecha de época lluviosa en el conjunto de
datos, y coincide con el pico más intenso del período. Con n=1 no se puede afirmar estacionalidad.
Lo correcto es decir que la única imagen de época lluviosa disponible muestra la floración más
intensa, lo cual es consistente con el mecanismo de lavado de nutrientes pero no lo demuestra.
Confirmar el patrón estacional requeriría un muestreo mensual balanceado durante varios años.}

% ═══════════════════════════════════════════════════════════════════════════════
\section{Conclusiones}
% ═══════════════════════════════════════════════════════════════════════════════

\begin{enumerate}[leftmargin=1.5em]
  \item \textbf{Amatitlán está en estado de eutrofia crónica.} Ninguna de las 11 fechas
        estudiadas muestra el lago limpio. El pico del 2026-06-19 (76.6\,\% del lago en
        floración densa, promedio de 60.2\mgm) es evidencia de un deterioro severo y urgente.

  \item \textbf{Atitlán mantiene buena calidad a escala de lago}, pero los focos puntuales
        en bahías cercanas a municipios con descarga de aguas residuales indican que la
        presión antrópica ya es perceptible desde el satélite. El análisis espacial es
        aquí la métrica clave; el promedio del lago enmascara estos eventos.

  \item \textbf{La diferencia entre ambos lagos no es de grado sino de régimen.} Amatitlán
        recibe en proporción mucho más carga de nutrientes respecto a su capacidad de dilución.
        Atitlán aún dispone de margen, pero ese margen no es ilimitado.

  \item \textbf{NDVI y NDWI refuerzan el diagnóstico en Amatitlán} (correlación significativa
        con Chl-a), pero son insuficientes en Atitlán porque la variación a escala de lago es
        demasiado pequeña para detectarla en el promedio. El análisis espacial a escala de píxel
        es imprescindible para Atitlán.

  \item \textbf{La teledetección no reemplaza el monitoreo en campo}, pero sí permite priorizar
        cuándo y dónde concentrar los recursos de muestreo físico, con una cobertura espacial y
        temporal que ningún equipo de campo puede igualar.
\end{enumerate}

% ═══════════════════════════════════════════════════════════════════════════════
\section*{Referencias}
\addcontentsline{toc}{section}{Referencias}
% ═══════════════════════════════════════════════════════════════════════════════

\begin{itemize}[leftmargin=1.5em]
  \item Kravitz, J. \& Matthews, M. (2020). \emph{Cyanobacteria Detection Script (cyanobacteria\_chla\_ndci\_l1c)}. Sentinel Hub Custom Scripts. \url{https://custom-scripts.sentinel-hub.com}

  \item Drusch, M. et al. (2012). Sentinel-2: ESA's Optical High-Resolution Mission for GMES Operational Services. \emph{Remote Sensing of Environment}, 120, 25--36.

  \item openEO platform. Copernicus Data Space Ecosystem. \url{https://openeo.dataspace.copernicus.eu}

  \item Autoridad para el Manejo Sustentable del Lago de Amatitlán (AMSA). Informes de calidad del agua. \url{https://www.amsa.gob.gt}

  \item Comité Nacional de Protección del Lago Atitlán (CONAP). Fichas de monitoreo ecológico.
\end{itemize}

\vfill
\begin{center}
  {\color{uvgblue}\rule{0.6\linewidth}{0.4pt}}\\[4pt]
  {\small Universidad del Valle de Guatemala · CC3084 Data Science · Semestre II 2026}
\end{center}

\end{document}
